% Options for packages loaded elsewhere
\PassOptionsToPackage{unicode}{hyperref}
\PassOptionsToPackage{hyphens}{url}
\PassOptionsToPackage{dvipsnames,svgnames,x11names}{xcolor}
\documentclass[
  10pt,
]{article}
\usepackage{xcolor}
\usepackage[margin=2.4cm]{geometry}
\usepackage{amsmath,amssymb}
\setcounter{secnumdepth}{-\maxdimen} % remove section numbering
\usepackage{iftex}
\ifPDFTeX
  \usepackage[T1]{fontenc}
  \usepackage[utf8]{inputenc}
  \usepackage{textcomp} % provide euro and other symbols
\else % if luatex or xetex
  \usepackage{unicode-math} % this also loads fontspec
  \defaultfontfeatures{Scale=MatchLowercase}
  \defaultfontfeatures[\rmfamily]{Ligatures=TeX,Scale=1}
\fi
\usepackage{lmodern}
\ifPDFTeX\else
  % xetex/luatex font selection
  \setmainfont[]{Libertinus Serif}
  \setmonofont[Scale=0.85]{Latin Modern Mono}
  \setmathfont[]{Latin Modern Math}
\fi
% Use upquote if available, for straight quotes in verbatim environments
\IfFileExists{upquote.sty}{\usepackage{upquote}}{}
\IfFileExists{microtype.sty}{% use microtype if available
  \usepackage[]{microtype}
  \UseMicrotypeSet[protrusion]{basicmath} % disable protrusion for tt fonts
}{}
\makeatletter
\@ifundefined{KOMAClassName}{% if non-KOMA class
  \IfFileExists{parskip.sty}{%
    \usepackage{parskip}
  }{% else
    \setlength{\parindent}{0pt}
    \setlength{\parskip}{6pt plus 2pt minus 1pt}}
}{% if KOMA class
  \KOMAoptions{parskip=half}}
\makeatother
\usepackage{longtable,booktabs,array}
\newcounter{none} % for unnumbered tables
\usepackage{calc} % for calculating minipage widths
% Correct order of tables after \paragraph or \subparagraph
\usepackage{etoolbox}
\makeatletter
\patchcmd\longtable{\par}{\if@noskipsec\mbox{}\fi\par}{}{}
\makeatother
% Allow footnotes in longtable head/foot
\IfFileExists{footnotehyper.sty}{\usepackage{footnotehyper}}{\usepackage{footnote}}
\makesavenoteenv{longtable}
\setlength{\emergencystretch}{3em} % prevent overfull lines
\providecommand{\tightlist}{%
  \setlength{\itemsep}{0pt}\setlength{\parskip}{0pt}}

\providecommand{\E}{\mathbb{E}}
\providecommand{\N}{\mathbb{N}}
\providecommand{\Z}{\mathbb{Z}}
\providecommand{\R}{\mathbb{R}}
\providecommand{\eps}{\varepsilon}
\providecommand{\ceil}[1]{\left\lceil #1\right\rceil}
\providecommand{\floor}[1]{\left\lfloor #1\right\rfloor}
\AtBeginDocument{\let\setminus\smallsetminus}
\usepackage[most]{tcolorbox}
\newtcolorbox{warning}{colback=red!5,colframe=red!65!black,boxrule=0.8pt,arc=2pt,left=6pt,right=6pt,top=5pt,bottom=5pt}
\usepackage{etoolbox}
\AtBeginEnvironment{longtable}{\small}
\setlength{\emergencystretch}{3em}
\usepackage{fancyhdr}
\pagestyle{fancy}
\fancyhf{}
\fancyhead[L]{\small\bfseries UNREFEREED GPT-6 ASTRA OUTPUT}
\fancyhead[R]{\small\thepage}
\renewcommand{\headrulewidth}{0.4pt}
\usepackage{bookmark}
\IfFileExists{xurl.sty}{\usepackage{xurl}}{} % add URL line breaks if available
\urlstyle{same}
\hypersetup{
  pdftitle={Three-chromatic (0{,}2)-graphs},
  pdfauthor={Writeup: gpt-6-astra (single model pass)},
  colorlinks=true,
  linkcolor={blue!50!black},
  filecolor={Maroon},
  citecolor={Blue},
  urlcolor={blue!50!black},
  pdfcreator={LaTeX via pandoc}}

\title{Three-chromatic (0,2)-graphs}
\usepackage{etoolbox}
\makeatletter
\providecommand{\subtitle}[1]{% add subtitle to \maketitle
  \apptocmd{\@title}{\par {\large #1 \par}}{}{}
}
\makeatother
\subtitle{Unrefereed candidate proof by GPT-6 Astra}
\author{Writeup: \texttt{gpt-6-astra} (single model pass)}
\date{Catalog id \texttt{three\_chromatic\_0\_2\_graphs} --- generated
2026-09-16}

\begin{document}
\maketitle

\begin{warning}

\textbf{UNREFEREED MODEL OUTPUT.} This document was selected solely
because GPT-6 Astra labelled its own result
\texttt{would\_publish:\ true} and returned \texttt{proved} or
\texttt{disproved}. It has not passed the adversarial LLM referee used
for the earlier Sol campaign, has not been checked by a human
mathematician, and has not been checked for novelty. Treat every
mathematical and bibliographic claim below as unverified.

\end{warning}

\section{Summary and provenance}\label{summary-and-provenance}

{\def\LTcaptype{none} % do not increment counter
\begin{longtable}[]{@{}
  >{\raggedright\arraybackslash}p{(\linewidth - 2\tabcolsep) * \real{0.5000}}
  >{\raggedright\arraybackslash}p{(\linewidth - 2\tabcolsep) * \real{0.5000}}@{}}
\toprule\noalign{}
\endhead
\bottomrule\noalign{}
\endlastfoot
Catalog id & \texttt{three\_chromatic\_0\_2\_graphs} \\
Catalog entry &
\href{https://graph-theory-ai.github.io/graph-conjectures/op/three_chromatic_0_2_graphs/}{Three-chromatic
(0,2)-graphs} \\
Source corpus & opg \\
Campaign leg & OpenProblemGarden first pass (\texttt{attacks\_opg}) \\
Source paper / entry &
http://www.openproblemgarden.org/op/three\_chromatic\_0\_2\_graphs \\
Model verdict & \textbf{proved} (confidence: high) \\
Model's one-line claim & Every finite 3-colourable (0,2)-graph is
bipartite, so none has chromatic number exactly three. \\
Model & \texttt{gpt-6-astra}, reasoning effort \texttt{max},
\texttt{mode=pro}, flex service tier \\
Original artifact &
\texttt{attacks\_opg/three\_chromatic\_0\_2\_graphs/output.md} \\
Independent review & \textbf{None} \\
Model's caveats & For finite simple undirected graphs; literature
priority has not been independently checked. \\
\end{longtable}
}

\section{Problem statement}\label{problem-statement}

Question Are there any (0,2)-graphs with chromatic number exactly three?

\subsection{Catalog context}\label{catalog-context}

A (0,2)-graph is a graph such that every pair of distinct vertices has
either 0 or 2 common neighbours. It is fairly easy to see that a
(0,2)-graph is necessarily regular and a variety of other properties can
be shown to hold. Although (0,2)-graphs with chromatic number 2, 4 and 5
are known it is open as to whether there can be any (0,2)-graphs with
chromatic number exactly three.

\section{Astra writeup}\label{astra-writeup}

\emph{The text below is the model output verbatim, apart from moving its
machine-readable verdict block into the summary above and shifting
Markdown heading levels for this document. Mechanical TeX defects and
equation tags were normalized where needed for compilation.}

\subsubsection{Statement}\label{statement}

\textbf{Theorem.} Let \(G\) be a finite simple graph in which every pair
of distinct vertices has either zero or two common neighbours. If \(G\)
admits a proper colouring with three colours, then \(G\) is bipartite.

The proof uses a neighbourhood-matching property to show that every
real-valued edge function whose sum around each quadrangle is zero is a
vertex-potential difference. A proper three-colouring supplies precisely
such an edge function.

\subsubsection{1. The neighbourhood-matching
property}\label{the-neighbourhood-matching-property}

For an edge \(uv\), put \[
A=N(u)\setminus\{v\},\qquad B=N(v)\setminus\{u\}.
\]

For each \(a\in A\), the vertices \(v,a\) have the common neighbour
\(u\), and hence have a unique second common neighbour \(b\). Define
\(\phi(a)=b\). Then \(b\in B\), \(ab\in E(G)\), and \[
u,v,b,a,u
\] is a quadrangle with four distinct vertices.

This map is a bijection. Indeed, \(u,b\) have \(v,a\) as their two
common neighbours. Consequently, the analogous construction from \(B\)
to \(A\)---taking the common neighbour of \(u,b\) other than \(v\)---is
its inverse.

Thus, for every edge \(uv\), there is a bijection \[
\phi:N(u)\setminus\{v\}\longrightarrow N(v)\setminus\{u\}
\quad\text{with}\quad a\phi(a)\in E(G).
\qquad\text{(1)}
\] In particular, adjacent vertices have equal degrees, so each
connected component is regular. The matching property itself, rather
than regularity, is what we need.

\subsubsection{2. A lemma about edge
functions}\label{a-lemma-about-edge-functions}

An \textbf{antisymmetric edge function} assigns a real number
\(\omega(x,y)\) to every ordered adjacent pair, with \[
\omega(y,x)=-\omega(x,y).
\] Say it is \textbf{square-closed} if its sum around every quadrangle
is zero.

\textbf{Lemma.} Suppose a finite connected graph satisfies (1). Every
square-closed antisymmetric edge function \(\omega\) has the form \[
\omega(x,y)=h(y)-h(x)
\qquad\text{(2)}
\] for some \(h:V(G)\to\mathbb R\).

\textbf{Proof.} Define \[
\operatorname{div}\omega(x)=\sum_{y\in N(x)}\omega(x,y),
\qquad
\Delta h(x)=\sum_{y\in N(x)}(h(y)-h(x)).
\] Antisymmetry gives \[
\sum_x\operatorname{div}\omega(x)=0.
\] The Laplacian of a finite connected graph has image equal to the
space of zero-sum functions. For completeness, it is symmetric, and \[
\langle h,-\Delta h\rangle
=\sum_{\{x,y\}\in E(G)}(h(x)-h(y))^2
\] shows that its kernel consists exactly of the constants. Hence there
exists \(h\) satisfying \[
\Delta h=\operatorname{div}\omega.
\]

Set \[
\eta(x,y)=\omega(x,y)-(h(y)-h(x)).
\] Then \(\eta\) is antisymmetric, square-closed, and divergence-free:
\[
\sum_{y\in N(x)}\eta(x,y)=0
\quad\text{for every }x.
\qquad\text{(3)}
\]

Suppose \(\eta\) is not identically zero. By finiteness and
antisymmetry, its maximum over ordered edges is some \(M>0\). Choose an
ordered edge \(uv\) with \[
\eta(u,v)=M.
\]

Use the matching (1). For \(a\in A\), write \(b=\phi(a)\).
Square-closedness on \(u,v,b,a,u\) gives \[
\eta(u,v)+\eta(v,b)-\eta(a,b)-\eta(u,a)=0.
\] Therefore \[
\eta(v,b)-\eta(u,a)=\eta(a,b)-M\le 0.
\] Summing over the matching yields \[
\sum_{b\in B}\eta(v,b)
\le
\sum_{a\in A}\eta(u,a).
\qquad\text{(4)}
\] But (3) and antisymmetry identify the two sides as \[
\sum_{b\in B}\eta(v,b)=-\eta(v,u)=M,
\qquad
\sum_{a\in A}\eta(u,a)=-\eta(u,v)=-M.
\] Thus (4) says \(M\le -M\), contradicting \(M>0\).

Consequently \(\eta=0\), proving (2). \(\square\)

\subsubsection{3. Applying the lemma to a
three-colouring}\label{applying-the-lemma-to-a-three-colouring}

Work in a connected component of \(G\); a one-vertex component is
harmless. Let \[
c:V(G)\longrightarrow \mathbb Z/3\mathbb Z
\] be a proper colouring. For each ordered edge \(xy\), define
\(\omega(x,y)\) to be the unique number in \(\{-1,1\}\) satisfying \[
\omega(x,y)\equiv c(y)-c(x)\pmod 3.
\qquad\text{(5)}
\] This is well-defined and antisymmetric.

Consider any quadrangle. The sum of \(\omega\) around it is:

\begin{itemize}
\tightlist
\item
  divisible by \(3\), by telescoping (5); and
\item
  an even integer between \(-4\) and \(4\), being a sum of four numbers
  from \(\{-1,1\}\).
\end{itemize}

The only possibility is zero. Thus \(\omega\) is square-closed.

The lemma now provides \(h:V(G)\to\mathbb R\) such that, on every edge,
\[
h(y)-h(x)=\omega(x,y)\in\{-1,1\}.
\qquad\text{(6)}
\] Subtract a constant so that \(h(r)=0\) at one chosen vertex.
Connectivity and (6) then imply that every \(h(x)\) is an integer.
Moreover, adjacent vertices have opposite parities. Hence \[
x\longmapsto h(x)\pmod 2
\] is a proper two-colouring.

Applying this independently to every component proves that \(G\) is
bipartite. Therefore \textbf{no finite \((0,2)\)-graph has chromatic
number exactly three}. \(\square\)

\subsubsection{Scope and gaps}\label{scope-and-gaps}

All finite simple graphs in the question are covered, including
disconnected graphs by the componentwise argument. No Cayley-graph or
cube-like hypothesis is used, and there is no remaining mathematical
case in the proof. I have not independently checked whether this
argument or the resulting theorem has already appeared in the
literature.

\end{document}
